\documentclass[11pt,a4paper]{article}
\usepackage[T1]{fontenc}
\usepackage{lmodern}
\usepackage{amsmath,amssymb,mathtools,mathrsfs}
\usepackage[margin=25mm,headheight=14pt]{geometry}
\usepackage{enumitem,booktabs,tabularx,fancyhdr,microtype}
\usepackage[hidelinks]{hyperref}
\hypersetup{pdftitle={An Analytic Hamiltonian KAM Theorem: Statement of the Current Lean Interface}}
\pagestyle{fancy}
\fancyhf{}
\fancyhead[L]{Analytic Hamiltonian KAM theorem}
\fancyhead[R]{Current Lean interface}
\fancyfoot[C]{\thepage}
\setlength{\parindent}{0pt}
\setlength{\parskip}{5pt}
\setlist[enumerate]{leftmargin=*,itemsep=5pt,topsep=4pt}
\newcommand{\R}{\mathbb R}
\newcommand{\C}{\mathbb C}
\newcommand{\Z}{\mathbb Z}
\newcommand{\T}{\mathbb T}
\newcommand{\norm}[1]{\lVert #1\rVert}
\newcommand{\code}[1]{\texttt{\small #1}}
\newcommand{\Gr}{G_{\R}}
\newcommand{\DD}{\mathcal D}
\newcommand{\PP}{\mathcal P}
\begin{document}
\begin{center}
{\LARGE\bfseries An Analytic Hamiltonian KAM Theorem}\par\vspace{5pt}
{\large Statement of the Current Lean Interface}\par\vspace{6pt}
Revised English formulation \quad | \quad October 8, 2026
\end{center}

This note states the mathematical input and output of
\code{KamProject.Arnold1963.theorem1}. It presents the project's basic
nondegenerate analytic Hamiltonian KAM theorem, developed with reference to
Arnold (1963), under the explicit assumptions of the current code. The source
version and changes from the earlier English draft are recorded on page~\pageref{sec:alignment}.

\section{Fixed data, domains, and conventions}

Fix an integer \(n\ge1\). Identify \(\R^n\) with its coordinatewise inclusion
in \(\C^n\), and use the coordinate maximum norm on both spaces. The norm on
\(\C^n\times\C^n\) is the product maximum norm. Coordinates are ordered as
\((p,q)\), with \(p\) the action and \(q\) the angle.

Let \(G\subset\C^n\) be compact and invariant under complex conjugation.
Define its real slice by
\[
 \Gr:=\{p\in\R^n:p\in G\}.
\]
Assume that a nonempty open set \(U\subset\R^n\) is specified such that
\begin{equation}\label{eq:geometry}
 \overline U=\Gr,\qquad
 m_n(\partial U)=0,\qquad
 U\subset\operatorname{int}_{\C^n}G,
\end{equation}
where closures and boundaries of real sets are taken in \(\R^n\), and
\(m_n\) is Lebesgue measure. In the last inclusion, real points are viewed
as complex points. These assumptions imply that \(U\) is bounded and that
\(0<m_n(\Gr)<\infty\). No complex-interior condition is imposed on every
point of \(\overline U\).

Let \(H_0:\C^n\to\C\) be holomorphic on a neighborhood of \(G\), and assume
\begin{equation}\label{eq:hzero}
 H_0(\overline p)=\overline{H_0(p)}\quad(p\in G),
 \qquad \det D^2H_0(p)\ne0\quad(p\in U).
\end{equation}
The Hessian uses complex coordinate derivatives. Its determinant is required
to be nonzero only on \(U\); the frequency map \(A=\nabla H_0\) is not
assumed globally injective.

Fix \(\rho>0\). Work on the complex covering space of the angles:
\[
 S_\rho:=\{q\in\C^n:\norm{\operatorname{Im}q}_\infty\le\rho\},
 \qquad \DD_\rho:=G\times S_\rho.
\]
The strip \(S_\rho\) has unbounded real directions. The physical phase space is
\[
 \T^n:=\R^n/(2\pi\Z)^n,\qquad
 \PP:=\Gr\times\T^n=\overline U\times\T^n.
\]
Write \(\pi(p,q)=(p,[q])\) for the real phase projection. On
\(\R^n\times\T^n\), let \(\mu\) be the product of action Lebesgue measure
and angle measure induced by a \(2\pi\)-periodic fundamental cell. Thus
\[
 0<\mu(\PP)=(2\pi)^n m_n(\Gr)=(2\pi)^n m_n(U)<\infty.
\]

\newpage
\section{Admissible perturbations and the main theorem}

A function \(H_1:\C^n\times\C^n\to\C\) is an admissible perturbation if:
\begin{enumerate}[label=\textup{(P\arabic*)}]
\item \textbf{Analyticity:} it is jointly holomorphic on an open neighborhood
of \(\DD_\rho\).
\item \textbf{Periodicity:} for every \((p,q)\in\DD_\rho\) and \(k\in\Z^n\),
\[
 H_1(p,q+2\pi k)=H_1(p,q).
\]
\item \textbf{Conjugation compatibility:} for every \((p,q)\in\DD_\rho\),
\[
 H_1(\overline p,\overline q)=\overline{H_1(p,q)}.
\]
In particular, \(H_0\) and \(H_1\) are real valued on their respective real
domains; this follows from the stated conjugation identities.
\item \textbf{Boundedness:} it is bounded on \(\DD_\rho\). Its domain norm is
\[
 \norm{H_1}_{G,\rho}:=\sup_{(p,q)\in\DD_\rho}|H_1(p,q)|<\infty.
\]
\end{enumerate}
The total functions above mirror the Lean representation. Analyticity and
the other domain properties are required only on the specified sets.

\subsection*{Theorem (the current \texttt{theorem1} interface)}
For fixed data satisfying Section~1, and for every \(\kappa>0\), there exists
\(M>0\), chosen before the perturbation, such that \emph{every} admissible
\(H_1\) with
\begin{equation}\label{eq:small}
 \norm{H_1}_{G,\rho}\le M
\end{equation}
has the following properties. Set
\[
 H(p,q):=H_0(p)+H_1(p,q),\qquad
 X_H(p,q):=(-\partial_qH(p,q),\partial_pH(p,q)).
\]
There exist \(r\) with \(0<r\le\rho\), sets \(F_{\rm good},F_{\rm bad}\subset\PP\),
and a family \(\mathscr T\) of actual subsets of \(\PP\), satisfying:
\begin{enumerate}[label=\textup{(C\arabic*)}]
\item \textbf{Partition and measure.} The good set is compact, both sets are
measurable, and
\[
 \PP=F_{\rm good}\,\dot\cup\,F_{\rm bad},\qquad
 F_{\rm good}\ne\varnothing,
\]
\[
 \mu(F_{\rm good})>0,\qquad
 (1-\kappa)_+\mu(\PP)<\mu(F_{\rm good}),\qquad
 \mu(F_{\rm bad})<\kappa\mu(\PP),
\]
where \(a_+=\max\{a,0\}\). This is the real-valued reading of the code's
\code{ENNReal.ofReal} inequality. For \(0<\kappa<1\), the lower bound is
\((1-\kappa)\mu(\PP)\).
\item \textbf{Family of tori.} Distinct members of \(\mathscr T\) are disjoint,
and \(F_{\rm good}=\bigcup_{T\in\mathscr T}T\). Every \(T\in\mathscr T\)
has the common-width torus realization specified in Section~3.
\item \textbf{Global trajectories.} For each \(z\in F_{\rm good}\), there is a
real-valued lifted trajectory \(\gamma:\R\to\C^n\times\C^n\) with
\(\pi(\gamma(0))=z\), \(\gamma'(t)=X_H(\gamma(t))\), and
\(\pi(\gamma(t))\in F_{\rm good}\) for every \(t\in\R\).
\end{enumerate}

\newpage
\section{Realization of each invariant torus}

For every \(T\in\mathscr T\), there exist a center \(c_T\in\Gr\), a frequency
\(\omega_T\in\R^n\), a lift \(L_T=(a_T,b_T):\C^n\to\C^n\times\C^n\),
an embedding \(e_T:\T^n\to\R^n\times\T^n\), and a map
\(J_T:\C^n\to\C^n\), with all of the following properties.

\begin{enumerate}[label=\textup{(T\arabic*)}]
\item \textbf{Frequency and center.}
\[
 \nabla H_0(c_T)=\omega_T,\qquad
 k\cdot\omega_T\ne0\quad(k\in\Z^n\setminus\{0\}).
\]
The family is indexed by its actual torus sets. Distinct tori are allowed
to have the same frequency.

\item \textbf{Joint angle analyticity, periodicity, and reality.}
The map \(L_T\) is holomorphic on a neighborhood of \(S_r\), and for
\(q\in S_r\), \(k\in\Z^n\),
\[
 L_T(q+2\pi k)=(a_T(q),b_T(q)+2\pi k).
\]
Moreover, \(L_T(q)\in\R^n\times\R^n\) for every \(q\in\R^n\).
Analyticity here is joint in all complex angle coordinates for each fixed
torus. No regularity in an open set of torus labels is asserted.

\item \textbf{Embedding and immersion.} The same lift induces the same torus:
\[
 e_T([q])=\pi(L_T(q))\quad(q\in\R^n),\qquad
 e_T(\T^n)=T.
\]
The map \(e_T\) is a closed topological embedding, and the complex
derivative \(DL_T(q)\) is injective at every \(q\in S_r\).
Thus the lift supplies an analytic embedded real \(n\)-torus.

\item \textbf{Angle inverse.} The angle component
\(Q\mapsto(e_T(Q))_q\) is a homeomorphism of \(\T^n\).
The map \(J_T\) is holomorphic near every real point, and
\[
 J_T(b_T(q))=q\quad(q\in\R^n),\qquad
 b_T(J_T(Q))=Q\quad(Q\in\R^n).
\]
These are the inverse identities recorded by the interface; no common
complex strip for \(J_T\) is required.

\item \textbf{Displacement.} For every \(q\in S_r\),
\[
 \norm{a_T(q)-c_T}_\infty<\kappa,\qquad
 \norm{b_T(q)-q}_\infty<\kappa.
\]
In particular, on real angles the parametrization has the familiar form
\(p=c_T+f_T(q)\), \([q_{\rm physical}]=[q+g_T(q)]\), with periodic
analytic corrections \(f_T=a_T-c_T\) and \(g_T=b_T-\mathrm{id}\).

\item \textbf{Original Hamiltonian dynamics for all time.} For every
\(q\in\R^n\), the real-valued curve
\[
 \gamma_{T,q}(t):=L_T(q+t\omega_T)\qquad(t\in\R)
\]
satisfies, for every real \(t\),
\[
 \frac{d}{dt}\gamma_{T,q}(t)=X_{H_0+H_1}(\gamma_{T,q}(t)),\qquad
 \pi(\gamma_{T,q}(t))\in T.
\]
Consequently the induced motion on \(T\) is conjugate via \(e_T\) to
\([q]\mapsto[q+t\omega_T]\). This expresses invariance and quasi-periodic
motion through the explicit orbit equation used by the code.
\end{enumerate}

\newpage
\section{Source alignment and changes from the earlier draft}\label{sec:alignment}

\textbf{Version used.} Local checkout: \code{publication/arnold1963-kam}.
Source commit:
\begin{center}
\code{f3c70ad77a0005331aca2a244f8f336677fdd454}.
\end{center}
Namespace: \code{KamProject.Arnold1963}.\par
Source paths below are relative to \code{KamProject/Arnold1963/}.

\begin{center}
\small
\begin{tabularx}{\linewidth}{@{}>{\raggedright\arraybackslash}p{0.30\linewidth}>{\raggedright\arraybackslash}X@{}}
\toprule
English specification & Lean declaration and source file\\
\midrule
Section 1 fixed data & \code{GlobalHamiltonianData} in
\code{Geometry/LocalizationCover.lean}\\[4pt]
Perturbations and norm & \code{AnalyticPhaseFunction}, \code{uniformNorm} in
\code{Basic/Functions.lean}\\[4pt]
Threshold and global output & \code{theorem1}, \code{Theorem1Result} in
\code{Main/Theorem1.lean}\\[4pt]
Individual torus data & \code{KAMTorus} in \code{Main/TorusResult.lean}\\[4pt]
Global family and disjointness & \code{FiniteLocalization.tori} and
\code{tori\_pairwise\_disjoint} in \code{Tori/CrossChart.lean}\\
\bottomrule
\end{tabularx}
\end{center}

\textbf{Changes to the mathematical statement.}
\begin{enumerate}[label=\arabic*.]
\item Replace the earlier complex-neighborhood condition on all of
\(\overline U\) by \eqref{eq:geometry}, including the exact real-slice
identity. Impose Hessian nondegeneracy on \(U\). Use
\(\overline U\times\T^n\) for the output partition and \(c_T\in\overline U\)
for the center condition.
\item State conjugation identities explicitly. Real-valuedness on the real
slice is a consequence, without an unsupported converse on other complex
components. Define the complex domain on the angle covering space, as in
\code{phaseDomain}, and include boundedness explicitly.
\item Put the quantifiers in the order \(\forall\kappa>0\;\exists M>0\;
\forall H_1\), and use \(\norm{H_1}_{G,\rho}\le M\). The threshold depends
only on the fixed unperturbed data, \(\rho\), and \(\kappa\). The companion
\code{theorem1\_of\_pointwise} accepts \(|H_1(z)|<M\) at every domain point;
its proof uses the resulting norm bound \(\norm{H_1}_{G,\rho}\le M\).
\item Replace the frequency-indexed ``foliation'' by a family of actual
torus sets. Retain the common positive width, compactness, positive measure,
both strict measure bounds, and the nonempty good set supplied by the code.
\item Express all torus properties using the same lift, embedding, center,
frequency, and angle inverse. Keep the displacement bound on \(S_r\) and
the orbit equation for the original Hamiltonian at every real time.
\end{enumerate}

\textbf{Relation to the literature.} This is a mathematical rendering of
the specified Lean interface. It does not assert equivalence with the full
range of hypotheses in Arnold's original paper. Its classical reference is
V. I. Arnol'd, \emph{Proof of a theorem of A. N. Kolmogorov on the invariance
of quasi-periodic motions under small perturbations of the Hamiltonian},
Russian Mathematical Surveys \textbf{18}(5) (1963), 9--36, especially
Section 2.1. This revision aligns the stated interface; it is not a new
audit of the entire formal proof.
\end{document}
